\documentclass[12pt, oneside, openany]{book}

\usepackage{mathptmx} % Contiene una fuente similar a Times New Roman

\usepackage[spanish, es-tabla]{babel} % Permite escritura en castellano
\usepackage[utf8]{inputenc} % Permite utilizar caracteres UTF8

\usepackage{graphicx} % Para la inclusión de gráficos e imágenes
\graphicspath{ {images/} } % Ruta para buscar las imágenes
\usepackage[a4paper,top=30mm,left=30mm,right=25mm,bottom=25mm,headheight=20mm]{geometry} % Configuración de los margenes de la página

% Paquetes para que funcione el formato.
\usepackage{titlesec}
\usepackage{setspace}
\usepackage{ragged2e}
\usepackage{fancyhdr}
\usepackage{lastpage}
\usepackage{stackengine}
\usepackage{array}
\usepackage[hidelinks]{hyperref}
\usepackage{enumitem}
\usepackage{float}
\usepackage{hypcap}
\usepackage{caption}
\usepackage{fancyvrb}

\usepackage{hyperref} % Paquete para que las referencias funcionen, y permite introducir links
\usepackage{xcolor} % Paquete para trabajar con colores (fondo de celdas, color del texto...)
% Se define un color gris desde su código RGB
\definecolor{gris}{RGB}{220,220,220}

\setcounter{secnumdepth}{3} % Para permitir numerar las sub-subsecciones

% Modifica el nombre de los índices al castellano
\addto\captionsspanish{
  \renewcommand{\contentsname}{Índice de contenido}
  \renewcommand{\listfigurename}{Índice de figuras}
  \renewcommand{\listtablename}{Índice de tablas}
}

% Formateo de los nombres de los apartados:
\titleformat{\chapter}[block]
  {\normalfont\Huge\bfseries\singlespacing}{\thechapter.}{1em}{\Huge}
\titlespacing*{\chapter}{0pt}{-62pt}{0pt}

\titleformat{\section}[block]
  {\normalfont\Large\bfseries}{\thesection.}{4pt}{\Large}
\titlespacing*{\section}{0pt}{\baselineskip}{0pt}

\titleformat{\subsection}[block]
  {\normalfont\large\bfseries}{\thesubsection.}{4pt}{\normalsize\large}
\titlespacing*{\subsection}{0pt}{0pt}{0pt}

\titleformat{\subsubsection}[block]
  {\normalfont\normalsize\bfseries}{\thesubsubsection.}{4pt}{\normalsize}
\titlespacing*{\subsubsection}{0pt}{0pt}{0pt}

\def\tablename{Tabla}

%% Variables para portada y cabeceras
%% Cambiar los valores para cada documento!!!
\def\title{Memoria de Prácticas}
\def\subject{Prácticas de Empresa}
\def\author{Juan Francisco Mier Montoto}
\def\target{www.mier.info}
\def\authorid{UO283319}
\def\date{enero 2024  }
\def\org{Escuela Politécnica de Ingeniería de Gijón}
\def\area{Grado en Ingeniería Informática en Tecnologías de la Información}

\def\ORG{\expandafter\MakeUppercase\expandafter{\org}}
\def\AREA{\expandafter\MakeUppercase\expandafter{\area}}
\def\SUBJECT{\expandafter\MakeUppercase\expandafter{\subject}}

\captionsetup{justification=centering}
\setlength{\headheight}{65pt}

\fancyhf{}
\fancyhead[L]{\includegraphics[height=16mm]{style/square.png}
  \hspace{1em} \Longstack[l] {
    \textbf{\SUBJECT} \newline
    \textbf{\title}}
  \newline \leftmark{}
}
\fancyhead[R]{\bfseries{Hoja \hyperlink{toc}{\thepage}~de~\pageref{LastPage}}}
\fancyfoot[C]{\href{\target}{\author}}
\renewcommand{\headrulewidth}{0pt} % default is 0pt
\renewcommand{\footrulewidth}{0.4pt} % default is 0

\fancypagestyle{plain}{%
  \fancyhf{}
  \fancyhead[L]{\includegraphics[height=16mm]{style/square.png}
    \hspace{1em} \Longstack[l]{
      \textbf{\SUBJECT} \newline
      \textbf{\title}}}
  \fancyhead[R]{\bfseries{Hoja \hyperlink{toc}{\thepage}~de~\pageref{LastPage}}}
  \fancyfoot[C]{\href{\target}{\author}}
  \renewcommand{\headrulewidth}{0pt} % default is 0pt
  \renewcommand{\footrulewidth}{0.4pt} % default is 0pt
}

\pagestyle{fancy}

\restylefloat{table}



\begin{document}

\rmfamily % Fuente tipo Romana

% Portada de la memoria
\begin{titlepage}
    \centering
    \bfseries {
        \null{}
        \vspace{0cm}
        \begin{table}[h]
            \centering
            \begin{tabular}{m{10cm} m{1cm} m{3cm}}
                \vspace{0.2cm}
                \includegraphics[width=86mm]{style/full.png} &  & \vspace{1.52mm} \includegraphics[width=23mm]{style/square.png} \\
            \end{tabular}
        \end{table}

        \vspace{3\baselineskip}

        \Large{\ORG{} \\ \vspace{3\baselineskip}}
        \large {
            \AREA{} \\ \vspace{3\baselineskip}
            \subject{} \\ \vspace{2\baselineskip}

            % TRABAJO FIN DE GRADO/MÁSTER Nº XXXXXXXXX \vspace{\baselineskip} \\
            \title{} \\ \vspace{3\baselineskip}

            \author{} \\
            \authorid{} \\
            % TUTOR/ES: \\
            % D. APELLIDO1 APELLIDO2, Nombre \\
            % D. APELLIDO1 APELLIDO2, Nombre \\  \vspace{\baselineskip}

            \vspace{2\baselineskip}
            FECHA:\@ \date{}
        }
    }
\end{titlepage}


% Índice de contenido
\addcontentsline{toc}{chapter}{Índice de contenido} % Añade la referencia al índice de contenido
\hypertarget{toc}{}
\tableofcontents
\newpage

% Índice de figuras
\addcontentsline{toc}{chapter}{Índice de figuras}  % Añade la referencia al índice de contenido
\hypertarget{lof}{}
\listoffigures

\justify{} % Texto justificado
\setlength{\parskip}{\baselineskip} % Separación entre párrafos de 1 linea
\onehalfspacing{}

%% El contenido de la memoria, dividido en capítulos:
\chapter{Introducción}
\section{Sobre la empresa}

\section{Sobre el departamento}

\section{Sobre el alumno}

\chapter{Tareas desempeñadas}
\section{Descripción}

\section{Relación con la titulación}

\chapter{Análisis posterior}
\section{Rol desempeñado}
Desde el comienzo de mi estancia en \emph{Okticket}, he desempeñado la labor normal
de un desarrollador de back-end especializado en I+D, teniendo acceso total a la
infraestructura y arquiectura de la empresa.

\section{Aprendizaje}
Durante las prácticas he desarrollado mucho conocimiento tanto a nivel técnico como
funcional de la empresa. Aunque mi rol haya sido más dedicado al I+D y al back-end,
he tocado diversos campos que me parecían muy interesantes, como dashboards, gestión
del cloud, gestión de infraestructuras, implementación de nuevas herramientas de
calidad de código, entre otros.

Además de tener la oportunidad de desarrollar e implementar nuevas características,
he tenido la oportunidad de entender cómo funciona una empresa de tecnología, cómo
se gestionan un proyecto, cómo se comunican los equipos y cómo se resuelven los
problemas y los compromisos. Esto me ha permitido aprender mucho sobre cómo se trabaja
en un entorno profesional, y me ha dado una visión más amplia de cómo se desarrollan
los proyectos en la vida real.

\chapter{Conclusiones}


%% Esto incluirá la bibliografía correctamente en nuestro trabajo
\newpage % En una nueva página
\addcontentsline{toc}{chapter}{Bibliografía} % Añade la referencia al índice de contenido

\bibliographystyle{ieeetr} % Define el estilo de la bibliografía
\bibliography{biblio} % Indica el archivo que contiene la colección de citas

\nocite{template}

\end{document}
